% -----------------------------
% 執筆にあたっての留意点
% -----------------------------
% 結果の章は実験・調査で得られたデータや分析結果について述べる章です．結果の章はデータや分析結果のみを淡々と記し，その解釈や結果から得られた知見は「考察」の章で述べるようにしてください．論文は「事実と意見を分けること」が重要となります．結果の章は実験・調査から得られた事実のみを述べるパートとし，意見は考察の章で述べるという棲み分けを行います．
%
% 結果は，測定項目・データ（=図や表）ごとに段落を分けて記述します．一般的には，各段落は，以下の3つの要素を順に並べて構成します．

% 1. 何の目的のために，何を調べたのかの説明（例：実験協力者がどの程度タスクに集中できたかを調べるため，タスク中のページ閲覧時間を調査した）
% 2. 得られた図・表の解説（例：図Xは，実験条件とページ閲覧時間の関係を示している）
% 3a. 図や表から読み取れる客観的事実（例：図Xによると，条件Aよりも条件Bのほうがページ閲覧時間が大きいことがうかがえる）
% 3b. 統計的な裏付け（例：分散分析を行った結果，条件Bと条件Aの間に統計的有意差があることが確認された）

% 前章で提示した仮説を漏れなく立証できるよう，必要となる実験・調査結果を記述してください．
%
%
% -----------------------------
% 補足（参考：how-to-xx/how-to-write-a-paper/1-学術論文の書き方.md 3.5節・5章，
%              how-to-xx/how-to-write-a-paper/2-文書の書き方.md 1章）
% -----------------------------
%
% ■ 都合の悪い結果も隠さずに書きます
% 仮説が支持されなかった結果や，提案手法が既存手法に負けた結果も，事実として報告してください．
% 隠して後から見つかるより，正直に報告した上で考察の章で誠実に議論する方が，査読では確実に好印象です．
%
% ■ 統計的な裏付けは，p値だけでは不十分です
% 有意差の有無に加えて，効果量と信頼区間も報告してください．「有意差があった」だけでは，
% その差が実質的に意味のある大きさなのかが読者に伝わりません．
% 検定を行う際は，その前提（正規性など）を満たしているかも確認しておきます．
%
% ■ 各段落の第1文で，その段落が何の結果の話なのかを述べます
% 上の構成の「1. 何の目的のために，何を調べたのか」がこれに当たります．
% 段落の先頭に「〜を調べるために，〜を測定した」と置けば，読者は図表を見る前に段落の主題を把握できます．
%
% ■ 図表の体裁
% * 図表はページ（2段組ならカラム）の上部に配置します
% * 本文で最初に参照するページか，その次のページに置きます
% * 図中の文字サイズは，本文とほぼ同じ大きさにします（縮小して貼り付けた図は，文字が潰れがちです）
% * すべての図表を本文中から参照します．本文から参照されない図表を置いてはいけません
% * 図表は，キャプションと軸ラベルだけを見て内容が理解できる状態を目指します

% ----------------------
% 結果
% ----------------------


\section{結果}